\chapter{Talente}

Talente sind die grundlegenden Eigenschaften und Fertigkeiten eines Charakters.
Sie definieren seine Stärken, Schwächen und Spezialisierungen im Spiel und
beeinflussen direkte Proben sowie den Einsatz von Fähigkeiten und Waffen.

\section{Grundlagen der Talente}

\begin{rulebox}{Talente}
	Ein \term{Talent}{talent} beschreibt die Ausbildung, Erfahrung oder Begabung
	eines Charakters in einem grundlegenden Kompetenzbereich.
\end{rulebox}

Talente werden verwendet für:
\begin{itemize}
	\item \termref[Proben]{proben}
	\item \termref[Kombinationswürfe]{kombinationswurf}
	\item \termref[Angriffe]{angriff} und \termref[Verteidigung]{verteidigung} 
	\item als Skalierung für  \termref[Fähigkeiten]{faehigkeiten}
\end{itemize}

Jedes Talent beginnt auf \textbf{Rang 1} und wird mit
\term{Talentpunkten}{talentpunkte} gesteigert.
Der Rang eines Talents bestimmt die verwendete Würfelgröße bei Proben.

\section{Bedeutung der einzelnen Talente}

Die folgenden Beschreibungen definieren den typischen Einsatzbereich eines Talents.
Im Zweifel entscheidet der Spielleiter.

\begin{description}
	\item[\textbf{ \term{Stärke}{staerke}}] Körperkraft, rohe Gewalt und physische Belastbarkeit.\\
	Anwendungen: Nahkampf, Lebenspunkte
	
	\item[\textbf{ \term{Gewandtheit}{gewandtheit}}] Schnelligkeit, Koordination, Reflexe und präzise Bewegung.\\
	Anwendungen: Ausweichen, Fernkampf, Initiative.
	
	\item[\textbf{ \term{Kampf}{kampf}}] Waffenführung, Kampfreflexe und Paraden.\\
	Anwendungen: Angriffe, Parieren, Kampftaktik, Duelle.
	
	\item[\textbf{ \term{Wahrnehmung}{wahrnehmung}}] Aufmerksamkeit, Sinne, Resistenzen gegen psychologische Angriffe\\
	Anwendungen: Fallen entdecken, Hinterhalte erkennen, Details wahrnehmen.
	
	\item[\textbf{ \term{Heimlichkeit}{heimnlichkeit}}] Leises Vorgehen, Verbergen und unauffälliges Handeln.\\
	Anwendungen: Schleichen, Verstecken/ Verkleidung, Spuren verwischen, unbemerkt agieren.
	
	\item[\textbf{ \term{Wissen}{wissen}}] Bildung, Recherche, Erinnerung, Sprache und Theorie.\\
	Anwendungen: Fakten abrufen, Texte deuten, Runen/Schriften erkennen, Zusammenhänge verstehen.
	
	\item[\textbf{ \term{Gesellschaft}{gesellschaft}}] Auftreten, soziale Interaktion, Manipulation und Empathie.\\
	Anwendungen: Überzeugen, Lügen erkennen, Verhandeln, Kontakte knüpfen.
	
	\item[\textbf{ \term{Wildnisleben}{wildnisleben}}] Überleben, Orientierung, Naturkunde und Umgang mit Tieren.\\
	Anwendungen: Spuren verfolgen, Nahrung finden, Lager bauen, Wetter/Natur einschätzen.
\end{description}

\section{Talent-Rang, Würfel und Talentproben}

Der Rang (1-6) eines Talents bestimmt die Würfelgröße, mit der Proben abgelegt werden.
Talentproben folgen den allgemeinen Regeln für Proben
(siehe Grundregeln).

\section{Kosten zur Steigerung von Talenten}

Die Kosten zur Steigerung eines Talents hängen von dessen Kategorie ab.
Die Kosten beziehen sich jeweils auf den neu erreichten Rang.

\begin{center}
	\begin{tabular}{lp{9cm}}
		\toprule
		\textbf{Kosten} & \textbf{Talente} \\
		\midrule
		2 Talentpunkte pro Rang & Stärke, Gewandtheit, Kampf \\
		1 Talentpunkt pro Rang & Wahrnehmung, Heimlichkeit, Wissen, Gesellschaft, Wildnisleben \\
		\bottomrule
	\end{tabular}
\end{center}


\section{Legendäre Effekte der Talente}

Ein Talent kann unter bestimmten Voraussetzungen, entsprechenden Slots und Meisterproben den \termref[legendären Rang (Rang 6)]{legendaerer-rang} erreichen.

\begin{rulebox}{Legendäre Effekte}
	Legendäre Effekte sind besondere, dauerhaft gewählte Fähigkeiten,
	die an ein Talent gebunden sind.
	Der Zugriff setzt voraus, dass das Talent Rang 6 erreicht hat.
\end{rulebox}

Beim Erreichen des legendären Rangs wählt der Spieler
\textbf{einen} der zwei verfügbaren Effekte des Talents.

Legendäre Effekte:
\begin{itemize}
	\item verändern oder erweitern die Wirkung eines Talents
	\item erfordern häufig den Einsatz von Marken
\end{itemize}

\subsection*{Übersicht: Legendäre Effekte der Talente}

\begin{center}
	\begin{tabular}{lp{6.5cm}p{6.5cm}}
		\toprule
		\textbf{Talent} & \textbf{Effekt A} & \textbf{Effekt B} \\
		\midrule
		Stärke &
		\textbf{Rohgewalt:} Überwindung physischer Hindernisse ist deutlich erleichtert
		(z.\,B. Springen, Klettern, Durchbrechen von Wänden). &
		\textbf{Selbsterhalt:} Die Regeneration von Lebenspunkten ist verdoppelt.
		Gib 1 Marke aus, damit die nächste Heilung nochmals doppelte Wirkung hat. \\
		
		Gewandtheit &
		\textbf{Schnelligkeit:} Erhält eine zusätzliche Bonusaktion. &
		\textbf{Initiativevorteil:} Erhält maximale Initiative im Kampf (wirkt nicht bei Mini-Boss- und Bossgegnern). \\
		
		Kampf &
		\textbf{Formationsbefehl:} Zu Kampfbeginn dürfen sich alle Gruppenmitglieder einmal
		frei neu positionieren. &
		\textbf{Kampfintuition:} Gib 1 Marke aus, um die nächste Aktion eines Gegners vorauszusehen. \\
		
		Wahrnehmung &
		\textbf{Rundumblick:} Beim Betreten eines Ortes erhältst du einen klaren Hinweis
		auf Besonderheiten. Fallen werden automatisch erkannt, sofern sie nicht übernatürlich verschleiert sind. &
		\textbf{Verhaltensanalyse:} Gib 1 Marke aus, um die Absicht oder Motivation einer Person zu erkennen. \\
		
		Heimlichkeit &
		\textbf{Spurenlos:} Der Charakter hinterlässt keine Spuren.
		Spuren der Gruppe können aktiv verwischt werden. &
		\textbf{Untertauchen:} Der Charakter und ein Begleiter können sich perfekt
		verstecken oder in Menschenmengen verschwinden. \\
		
		Wissen &
		\textbf{Sprachkundig}: Gängige Dialekte, Schriften und Codes können ohne Probe
		grundlegend verstanden werden. Bei Rätseln erhältst du pro Rätsel
		einen starken Hinweis. &
		\textbf{Genie:} Gib 1 Marke aus, um vom Spielleiter die bestmögliche Option
		oder die Bestätigung einer plausiblen Annahme zu erhalten. \\
		
		Gesellschaft &
		\textbf{Netzwerk:} In bewohnten Orten findest du einmal pro Ort und Tag
		einen hilfreichen Kontakt für kleine Gefälligkeiten. &
		\textbf{Überzeugend:} Gib 1 Marke aus, um von einer Zielperson eine kleine Wahrheit
		oder ein zurückgehaltenes Detail zu erfahren. \\
		
		Wildnisleben &
		\textbf{Kräuterkundig:} Findet überall nützliche Kräuter (siehe Ausrüstung). &
		\textbf{Tierfreund:} Neutrale Tiere versuchen dir im Rahmen ihrer Möglichkeiten
		zu helfen. \\
		\bottomrule
	\end{tabular}
\end{center}



